\section{Wnioski}

Celem pracy była integracja, zaprojektowanie i implementacja komponentu renderującego obiekty 3d z częścią webową, w tym przypadku, systemem aukcyjnym.

Pierwszą częścią pracy nad systemem, było stworzenie projektu. Projekt zawierał wymagania funkcjonalne, niefunkcjonalne oraz diagramy ERD i BPMN. Wymienienie wymagań funkcjonalnych i niefunkcjonalnych pokazało zakres działania systemu. Diagram ERD ułatwił proces implementacji bazy danych, a diagramy BPMN pozwoliły na szybkie zrozumienie i wdrożenie procesów biznesowych w systemie aukcyjnym.

Kolejnym etapem pracy była implementacja komponentu renderującego. Proces ten wymagał dogłębnego zrozumienia niskopoziomowej komunikacji z procesorem graficznym poprzez API WebGPU. Zastosowanie wzorca potoku opartego na \texttt{std::expected} ze standardu C++23 pozwoliło na elegancką obsługę błędów podczas inicjalizacji trzynastu wymaganych zasobów GPU. Implementacja shaderów w języku WGSL umożliwiła realizację zaawansowanego modelu oświetlenia z mapami normalnych oraz wieloma źródłami światła. Szczególnym wyzwaniem okazała się kompilacja kodu do formatu WebAssembly przy użyciu narzędzia Emscripten, co wymagało odpowiedniej konfiguracji flag kompilacji oraz zrozumienia ograniczeń wirtualnego systemu plików przeglądarki. Dzięki temu udało się osiągnąć wydajność renderowania porównywalną z aplikacjami natywnymi, zachowując jednocześnie pełną kompatybilność z ekosystemem internetowym.

Następnie pracowaliśmy nad implementacją systemu aukcyjnego. System ten został podzielony na część internetową oraz serwerową. Część internetowa została stworzona zgodnie z początkowym projektem wizualnym. Podczas tworzenia części serwerowej bardzo pomogły diagramy BPMN i ERD, dzięki którym mogliśmy wdrożyć zaprojektowane rozwiązania.

Końcowym etapem było sprawdzenie działania aplikacji. Każda z funkcjonalności została przetestowana manualnie. Wykryte w ten sposób niepoprawne działania zostały naprawiane na bieżąco.

Cel pracy został w pełni zrealizowany. Udało nam się utworzyć w pełni działający kod, który pozwala na renderowanie obiektów 3D i opakowaliśmy go w system aukcyjny. Dzięki temu udało nam się znaleźć nowe zastosowanie renderingu w WebGPU.

Projektowanie tego systemu umożliwiło członkom grupy zdobycie praktycznego doświadczenia w pracy nad rozbudowaną aplikacją. Proces ten ukazał, jak istotna jest współpraca z wykorzystaniem zdalnego repozytorium oraz jak ważne są praktyki pisania czystego, czytelnego kodu.